\section{Report and monochrome table presets}
\key{style=matrix} provides centered cells/headers, row rules, row stretch 1.35,
column separation 6pt and row padding 1pt, with an independent color theme.
\begin{demo}{Matrix style with missing values and a summary}
\begin{center}
\begin{gradtable}[align=c,header align=c,style=matrix,theme=teal,width=\linewidth]
  \gradcolumn{name}{Method}
  \gradcolumn[type=number,aggregate=mean]{value}{Value}
  \gradrow{A,-5}\gradrow{B,5}\gradrow{Pending,NA}\gradsummary{Mean}
\end{gradtable}
\end{center}
\end{demo}
The optional \key{style=reference} uses booktabs rules, centered headers,
decimal number alignment, \key{row stretch=1.25} and \key{column sep=6pt}.
It has no mandatory color theme and leaves existing styles unchanged.
\begin{demo}{Reference table with aligned numbers}
\begin{center}
\begin{gradtable}[align=c,header align=c,style=reference,theme=teal,width=\linewidth]
  \gradcolumn{name}{Method}
  \gradcolumn[type=number,aggregate=mean,options={precision=3}]{value}{Value}
  \gradrow{A,-12.125}\gradrow{B,1.25}\gradrow{Pending,NA}\gradsummary{Mean}
\end{gradtable}
\end{center}
\end{demo}
Table presets are paper (thin solid bars and compact booktabs rules), report
(tinted headings and taller rows) and mono (gray and textures). Column drawing
options override the preset. Table presets and individual-graphic presets are
separate interfaces.
\begin{demo}{Bullet and interval columns}
\begin{center}
\begin{gradtable}[align=c,header align=c,preset=report,graphic width=32mm,
  text width=19mm,label width=19mm]
  \gradcolumn{name}{Project}
  \gradcolumn[type=bullet,options={max=100,precision=0,
    bullet bands={{60/black!8},{85/black!16},{100/black!25}},
    target=85}]{actual}{Progress}
  \gradcolumn[type=interval,options={max=100,precision=0,
    color=gradbarsTeal}]{estimate}{Interval}
  \gradrow{A,78,65/78/90}
  \gradrow{B,92,80/92/98}
  \gradrow{C,NA,NA}
\end{gradtable}
\end{center}
\end{demo}
Interval cells use lower/point/upper. NA in the point position omits only the
point, while whole-cell NA omits the interval. The supplied point must lie inside
the original endpoints. Bullet columns require explicit background bands.
\begin{demo}{Signed observations in monochrome}
\begin{center}
\begin{gradtable}[align=c,header align=c,preset=mono,graphic width=32mm,
  text width=19mm,label width=19mm]
  \gradcolumn{name}{Group}
  \gradcolumn[type=bar,options={min=-40,max=60,precision=0}]{change}{Change}
  \gradcolumn[type=interval,options={min=-40,max=60,
    precision=0,marker shape=diamond}]{range}{Interval}
  \gradrow{A,-12,-25/-12/5}
  \gradrow{B,28,15/28/45}
\end{gradtable}
\end{center}
\end{demo}
Default tables are indivisible and can be placed in table, table* or minipage.
They do not infer ranges, calculate totals or load files automatically.

\clearpage

\subsection{Built-in table styles and independent themes}
Style=academic uses paper rules, compact uses smaller text and natural widths,
and dashboard uses report spacing and row rules. They use the existing named
style registry; place a style before explicit overrides.
Table theme accepts none (default), blue, teal, violet, orange, forest or mono.
Themes set graphic fills, sequential heat endpoints and automatic header/summary
fills. Column drawing options can override colors or select a diverging palette.
\begin{demo}{Academic with violet}
\begin{center}
\begin{gradtable}[align=c,header align=c,style=academic,theme=violet,width=\linewidth]
  \gradcolumn{name}{Method}
  \gradcolumn[type=bar,aggregate=mean]{score}{Score}
  \gradrow{A,72}\gradrow{B,90}\gradsummary{Mean}
\end{gradtable}
\end{center}
\end{demo}
\begin{demo}{Compact with forest green}
\begin{center}
\begin{gradtable}[align=c,header align=c,style=compact,theme=forest,width=\linewidth]
  \gradcolumn{name}{Method}
  \gradcolumn[type=heat,aggregate=mean]{score}{Score}
  \gradrow{A,72}\gradrow{B,90}\gradsummary{Mean}
\end{gradtable}
\end{center}
\end{demo}
\begin{demo}{Dashboard with warm orange}
\begin{center}
\begin{gradtable}[align=c,header align=c,style=dashboard,theme=orange,width=\linewidth]
  \gradcolumn{name}{Method}
  \gradcolumn[type=bar,aggregate=mean]{score}{Score}
  \gradrow{A,72}\gradrow{B,90}\gradsummary{Mean}
\end{gradtable}
\end{center}
\end{demo}
Explicit header/summary colors override auto fills; empty colors disable them.
Font controls table text, while column options can style graphic label fonts.
Existing paper/report/mono presets remain available. Styles do not change
statistical scopes, zero baselines or numerical scales.


\subsection{The research table style}
Research uses booktabs, centered headers, row stretch 1.3, column separation
5pt, row padding 1pt and group separation 5pt. Number columns default to decimal
alignment. No color theme or SI dependency is mandatory; all six existing table
themes can be combined. Academic, compact and dashboard retain their defaults.
Put explicit overrides after style selection.
\begin{demo}{Research style with the forest theme}
\begin{center}
\begin{gradtable}[align=c,header align=c,style=research,theme=forest,width=\linewidth,group summary=true]
  \gradcolumn{name}{Method}
  \gradcolumn[type=number,aggregate=mean,mark=both,
    options={precision=3}]{value}{Value}
  \gradgroup{Experiment}\gradrow{A,-12.125}\gradrow{B,1.25}\gradrow{Pending,NA}
  \gradsummary{Overall}
\end{gradtable}
\end{center}
\end{demo}
